\documentclass[11pt, a4paper]{article}
\usepackage[a4paper, top=2.5cm, bottom=2.5cm, left=2cm, right=2cm]{geometry}
\usepackage{fontspec}
\usepackage[english, bidi=basic, provide=*]{babel}
\babelprovide[import, onchar=ids fonts]{english}
\babelfont{rm}{TeX Gyre Termes}
\babelfont{sf}{TeX Gyre Heros}
\babelfont{tt}{TeX Gyre Cursor}
\usepackage{enumitem}
\usepackage{amsmath}

\title{\textbf{Chapter 4: Carbon and the Molecular Diversity of Life \\ Exam Notes}}
\author{}
\date{}

\begin{document}

\maketitle

\section*{4.1 Organic Chemistry is Key to the Origin of Life}
\begin{itemize}
    \item \textbf{Organic Chemistry:} The study of carbon compounds, regardless of their origin. Most organic compounds contain hydrogen atoms in addition to carbon atoms.
    \item \textbf{Vitalism vs. Mechanism:} 
    \begin{itemize}
        \item \textit{Vitalism:} The historical, disproved belief that organic compounds could only arise in living organisms.
        \item \textit{Mechanism:} The accepted view that physical and chemical laws govern all natural phenomena. 
    \end{itemize}
    \item \textbf{Miller-Urey Experiment (1953):} Stanley Miller demonstrated the abiotic (non-living) synthesis of organic compounds. By passing an electrical spark through a mixture of gases ($H_2$, $CH_4$, $NH_3$, and water vapor) simulating the early Earth's atmosphere, he produced amino acids and other organic molecules. This supported the idea that complex organic molecules could arise spontaneously under early Earth conditions.
\end{itemize}

\section*{4.2 Carbon's Tetravalence Enables Molecular Diversity}
Carbon is unparalleled in its ability to form large, complex, and diverse molecules.
\begin{itemize}
    \item \textbf{Valence of Carbon:} Carbon has an atomic number of 6 (2 electrons in the first shell, 4 in the valence shell). To complete its octet, it forms \textbf{four covalent bonds} (tetravalence).
    \item \textbf{Molecular Shape:} 
    \begin{itemize}
        \item When a carbon atom forms four single covalent bonds, the arrangement of its four hybrid orbitals causes the bonds to angle toward the corners of an imaginary \textbf{tetrahedron} (bond angles of $109.5^\circ$). Example: Methane ($CH_4$).
        \item When two carbon atoms are joined by a \textbf{double bond}, the atoms joined to the carbons are in the same plane as the carbons (a flat, planar shape). Example: Ethene ($C_2H_4$).
    \end{itemize}
    \item \textbf{Carbon Skeletons:} Carbon chains form the skeletons of most organic molecules. Skeletons vary in four primary ways:
    \begin{enumerate}
        \item \textit{Length} (number of carbons).
        \item \textit{Branching} (straight chains vs. branched chains).
        \item \textit{Double bond position} (location and number of double bonds).
        \item \textit{Presence of rings} (e.g., cyclohexane, benzene).
    \end{enumerate}
    \item \textbf{Hydrocarbons:} Organic molecules consisting \textit{only} of carbon and hydrogen. They are \textbf{nonpolar} and \textbf{hydrophobic}. They can undergo reactions that release a large amount of energy (e.g., fossil fuels, the hydrocarbon tails of fats).
\end{itemize}

\subsection*{Isomers}
Isomers are compounds that have the same number of atoms of the same elements (same molecular formula) but different structures, and hence different properties.
\begin{itemize}
    \item \textbf{1. Structural Isomers:} Differ in the covalent arrangements of their atoms (e.g., straight-chain pentane vs. branched isopentane).
    \item \textbf{2. Cis-Trans Isomers (Geometric Isomers):} Carbons have covalent bonds to the same atoms, but these atoms differ in their spatial arrangements due to the \textbf{inflexibility of double bonds}.
    \begin{itemize}
        \item \textit{Cis isomer:} The two "X" groups are on the \textit{same} side of the double bond.
        \item \textit{Trans isomer:} The two "X" groups are on \textit{opposite} sides of the double bond.
    \end{itemize}
    \item \textbf{3. Enantiomers:} Isomers that are \textbf{mirror images} of each other and differ in shape due to the presence of an \textbf{asymmetric carbon} (a carbon bonded to four \textit{different} atoms or groups). 
    \begin{itemize}
        \item Usually, only one enantiomer is biologically active because only that form can bind to specific molecules in an organism.
        \item \textit{Example:} L-Dopa is effective against Parkinson's disease; D-Dopa is inactive. Ibuprofen and Albuterol also exist as enantiomers where one form is far more effective.
    \end{itemize}
\end{itemize}

\section*{4.3 Chemical Groups and Molecular Function}
The distinctive properties of an organic molecule depend not only on the arrangement of its carbon skeleton but also on the chemical groups attached to that skeleton.

\subsection*{The Seven Most Important Functional Groups}
Functional groups are directly involved in chemical reactions. (The methyl group is an exception; it serves as a biological tag rather than participating directly in reactions).
\begin{enumerate}
    \item \textbf{Hydroxyl Group ($-OH$):} 
    \begin{itemize}
        \item Is polar due to electronegative oxygen. Forms hydrogen bonds with water, helping dissolve compounds such as sugars.
        \item Compound name: \textbf{Alcohol} (specific name usually ends in \textit{-ol}, e.g., Ethanol).
    \end{itemize}
    \item \textbf{Carbonyl Group ($>C=O$):} 
    \begin{itemize}
        \item Sugars with ketone groups are called ketoses; those with aldehydes are called aldoses.
        \item \textbf{Ketone:} The carbonyl group is \textit{within} a carbon skeleton.
        \item \textbf{Aldehyde:} The carbonyl group is at the \textit{end} of a carbon skeleton.
    \end{itemize}
    \item \textbf{Carboxyl Group ($-COOH$):}
    \begin{itemize}
        \item Acts as an \textbf{acid} (can donate $H^+$) because the covalent bond between oxygen and hydrogen is so polar.
        \item Compound name: \textbf{Carboxylic acid} or \textbf{Organic acid} (e.g., Acetic acid).
    \end{itemize}
    \item \textbf{Amino Group ($-NH_2$):}
    \begin{itemize}
        \item Acts as a \textbf{base}; can pick up an $H^+$ from the surrounding solution (water, in living organisms).
        \item Compound name: \textbf{Amine} (e.g., Glycine, which has both a carboxyl and an amino group, making it an amino acid).
    \end{itemize}
    \item \textbf{Sulfhydryl Group ($-SH$):}
    \begin{itemize}
        \item Two $-SH$ groups can react, forming a "cross-link" that helps stabilize protein structure (e.g., hair protein).
        \item Compound name: \textbf{Thiol} (e.g., Cysteine).
    \end{itemize}
    \item \textbf{Phosphate Group ($-OPO_3^{2-}$):}
    \begin{itemize}
        \item Contributes negative charge (1- when positioned inside a chain of phosphates; 2- when at the end). When attached, confers on a molecule the ability to react with water, releasing energy.
        \item Compound name: \textbf{Organic phosphate} (e.g., Glycerol phosphate).
    \end{itemize}
    \item \textbf{Methyl Group ($-CH_3$):}
    \begin{itemize}
        \item \textit{Not} reactive, but affects the expression of genes when on DNA or on proteins bound to DNA. Affects the shape and function of male and female sex hormones.
        \item Compound name: \textbf{Methylated compound} (e.g., 5-Methylcytosine).
    \end{itemize}
\end{enumerate}

\subsection*{ATP: An Important Source of Energy for Cellular Processes}
\begin{itemize}
    \item \textbf{ATP (Adenosine Triphosphate):} Consists of an organic molecule called adenosine attached to a string of \textbf{three phosphate groups}.
    \item \textbf{Function:} Stores the potential to react with water. This reaction releases energy that can be used by the cell ($ATP + H_2O \rightarrow ADP + P_i + Energy$).
\end{itemize}

\end{document}
